\subsection{Direct-sum sector conjugations}

\begin{definition}[Direct-sum unitary conjugations]
    \label{def:mpdo_bnt_direct_sum_unitary_conjugations}
    \lean{MPOTensor.BNTAlgebraTensorClause.TwoSiteExactSectorGauge.%
        RelabeledTwoSiteSectorAlgebra}
    \lean{MPOTensor.BNTAlgebraTensorClause.TwoSiteExactSectorGauge.%
        UnitarySectorConjugacy}
    \lean{MPOTensor.BNTAlgebraTensorClause.TwoSiteExactSectorGauge.%
        UnitarySectorConjugacy.sectorLinearEquiv}
    \lean{MPOTensor.BNTAlgebraTensorClause.TwoSiteExactSectorGauge.%
        UnitarySectorConjugacy.directSumUnitaryT}
    \lean{MPOTensor.BNTAlgebraTensorClause.TwoSiteExactSectorGauge.%
        UnitarySectorConjugacy.directSumUnitaryS}
    \lean{MPOTensor.BNTAlgebraTensorClause.TwoSiteExactSectorGauge.%
        UnitarySectorConjugacy.directSumUnitaryT_apply}
    \lean{MPOTensor.BNTAlgebraTensorClause.TwoSiteExactSectorGauge.%
        UnitarySectorConjugacy.directSumUnitaryS_apply}
    \leanok
    \uses{def:mpdo_bnt_two_site_exact_sector_gauge}
    Work in the setting of
    Definition~\ref{def:mpdo_bnt_two_site_exact_sector_gauge}.  Let
    \begin{align}
        \mathcal A_1
        &=\bigoplus_\gamma \MN{d_\gamma},
        &
        \mathcal A_2^\sigma
        &=\bigoplus_\gamma \MN{\widetilde d_{\sigma(\gamma)}}
    \end{align}
    be the one-site sector algebra and the two-site sector algebra relabelled
    by the given sector matching.  Suppose that
    $d_\gamma=\widetilde d_{\sigma(\gamma)}$ and that there is a unitary
    $U_\gamma$ on each matched sector such that
    \begin{align}
        A_{\sigma(\gamma)}^i
        &=U_\gamma M_\gamma^i U_\gamma^*
        \label{eq:mpdo_bnt_sector_unitary_conjugacy}
    \end{align}
    for every physical index $i$.  Denote by $\iota_\gamma$ the coordinate
    identification induced by the equality of the two bond dimensions.
    Define mutually inverse linear maps
    \begin{align}
        \widetilde T\colon \mathcal A_1&\longrightarrow\mathcal A_2^\sigma,
        &
        [\widetilde T(X)]_\gamma
        &=U_\gamma\iota_\gamma(X_\gamma)U_\gamma^*,
        \label{eq:mpdo_bnt_direct_sum_unitary_T}
        \\
        \widetilde S\colon \mathcal A_2^\sigma&\longrightarrow\mathcal A_1,
        &
        [\widetilde S(Y)]_\gamma
        &=\iota_\gamma^{-1}
          (U_\gamma^*Y_\gamma U_\gamma).
        \label{eq:mpdo_bnt_direct_sum_unitary_S}
    \end{align}
    These are the converse-direction middle conjugations introduced in
    \cite[Appendix~C.4, lines~2053--2080]{Cirac2016MPDO_arXiv}.  The
    existence of the unitaries in
    (\ref{eq:mpdo_bnt_sector_unitary_conjugacy}) is an assumption here.
    They are distinct from the transported direct-sum maps in
    \cite[Appendix~C.4, lines~1955--1997]{Cirac2016MPDO_arXiv}.  No maps
    between the ambient full matrix algebras are included here.
\end{definition}

\begin{theorem}[Properties of the direct-sum unitary conjugations]
    \label{thm:mpdo_bnt_direct_sum_unitary_conjugations}
    \lean{MPOTensor.BNTAlgebraTensorClause.TwoSiteExactSectorGauge.%
        UnitarySectorConjugacy.directSumUnitaryT_isKrausDirectSumMap}
    \lean{MPOTensor.BNTAlgebraTensorClause.TwoSiteExactSectorGauge.%
        UnitarySectorConjugacy.directSumUnitaryS_isKrausDirectSumMap}
    \lean{MPOTensor.BNTAlgebraTensorClause.TwoSiteExactSectorGauge.%
        UnitarySectorConjugacy.%
        directSumUnitaryT_isTracePreservingBetweenDirectSums}
    \lean{MPOTensor.BNTAlgebraTensorClause.TwoSiteExactSectorGauge.%
        UnitarySectorConjugacy.%
        directSumUnitaryS_isTracePreservingBetweenDirectSums}
    \lean{MPOTensor.BNTAlgebraTensorClause.TwoSiteExactSectorGauge.%
        UnitarySectorConjugacy.directSumUnitaryS_comp_directSumUnitaryT}
    \lean{MPOTensor.BNTAlgebraTensorClause.TwoSiteExactSectorGauge.%
        UnitarySectorConjugacy.directSumUnitaryT_comp_directSumUnitaryS}
    \lean{MPOTensor.BNTAlgebraTensorClause.TwoSiteExactSectorGauge.%
        UnitarySectorConjugacy.directSumUnitaryT_tensor}
    \lean{MPOTensor.BNTAlgebraTensorClause.TwoSiteExactSectorGauge.%
        UnitarySectorConjugacy.directSumUnitaryS_tensor}
    \leanok
    \uses{def:kraus_map_between_direct_sums,
        def:mpdo_bnt_direct_sum_unitary_conjugations}
    The maps $\widetilde T$ and $\widetilde S$ are completely positive and
    preserve the total trace on the direct sums.  They are mutually inverse:
    \begin{align}
        \widetilde S\widetilde T
        &=\id_{\mathcal A_1},
        &
        \widetilde T\widetilde S
        &=\id_{\mathcal A_2^\sigma}.
        \label{eq:mpdo_bnt_direct_sum_unitary_inverses}
    \end{align}
    They also carry the matched tensor letters in both directions:
    \begin{align}
        \widetilde T\bigl((M_\gamma^i)_\gamma\bigr)
        &=\bigl(A_{\sigma(\gamma)}^i\bigr)_\gamma,
        \notag
        \\
        \widetilde S\bigl((A_{\sigma(\gamma)}^i)_\gamma\bigr)
        &=\bigl(M_\gamma^i\bigr)_\gamma.
        \label{eq:mpdo_bnt_direct_sum_unitary_tensor_letters}
    \end{align}
\end{theorem}

\begin{proof}\leanok
    \uses{lem:direct_sum_kraus_composition_schwarz}
    On each summand, conjugation by $U_\gamma$ and the coordinate
    identification $\iota_\gamma$ are completely positive and
    trace-preserving.  Their direct sums therefore have the same properties.
    Unitarity gives
    \begin{align}
        U_\gamma^*U_\gamma
        &=I,
        &
        U_\gamma U_\gamma^*
        &=I,
    \end{align}
    which proves (\ref{eq:mpdo_bnt_direct_sum_unitary_inverses}) sector by
    sector.  The first identity in
    (\ref{eq:mpdo_bnt_direct_sum_unitary_tensor_letters}) is precisely
    (\ref{eq:mpdo_bnt_sector_unitary_conjugacy}); the second follows by
    applying the inverse map.
\end{proof}
